% Unit 3 exam -- double block. 30 MCQ (no calc) + 1 long + 1 short FRQ.
\documentclass{apchem-exam}

\apcunit{3}
\apcunittitle{Properties of Substances and Mixtures}
\apcdoctitle{Unit Exam}
\apcsubtitle{Wed Sep 30 \textbullet\ double block, 76 minutes \textbullet\ periodic table and equations sheet provided}

\begin{document}
\apctitleblock

\examsection{I}{Multiple Choice}{30 questions \textbullet\ 45 minutes \textbullet\ 1 point each}{\nocalculator}

% ---- 3.1 IMFs (5) ----------------------------------------------------------
\begin{mcq}
  Which intermolecular force is present in \emph{all} molecular substances?
  \choices
    \choice{dipole--dipole}
    \choice{hydrogen bonding}
    \correct{London dispersion}
    \choice{ion--dipole}
\end{mcq}
\why{Dispersion arises from instantaneous dipoles in any electron cloud.}

\begin{mcq}
  Which substance exhibits hydrogen bonding?
  \choices
    \choice{\ce{CH3F}}
    \correct{\ce{CH3OH}}
    \choice{\ce{CH3OCH3}}
    \choice{\ce{H2S}}
\end{mcq}
\why{Requires H bonded directly to N, O, or F --- only methanol's O--H
qualifies.}
\distractor{C}{has O but no H attached to it}
\distractor{A}{F is present but bonded to C, not H}

\begin{mcq}
  \ce{BI3} (nonpolar) boils at \SI{210}{\celsius} while \ce{NH3}
  (hydrogen bonding) boils at \SI{-33}{\celsius}. The best explanation is
  that
  \choices
    \correct{\ce{BI3} has far more electrons, giving very strong dispersion
      forces}
    \choice{\ce{NH3} does not actually hydrogen bond}
    \choice{\ce{BI3} is ionic}
    \choice{boiling point does not depend on intermolecular forces}
\end{mcq}
\why{Force magnitude, not force type, sets the ranking when sizes differ
greatly.}

\begin{mcq}
  Which property \emph{decreases} as intermolecular forces get stronger?
  \choices
    \choice{boiling point}
    \choice{viscosity}
    \correct{vapor pressure}
    \choice{surface tension}
\end{mcq}
\why{Strong IMFs keep molecules in the liquid, so fewer escape to the vapor.}

% ---- 3.2-3.3 solids and phases (4) -----------------------------------------
\begin{mcq}
  A solid melts at \SI{1600}{\celsius}, is very hard, and does not conduct
  electricity in any phase. It is best classified as
  \choices
    \choice{ionic}
    \choice{metallic}
    \correct{covalent network}
    \choice{molecular}
\end{mcq}
\why{Covalent bonds throughout the lattice with no mobile charges.}
\distractor{A}{ionic solids conduct when molten}

\begin{mcq}
  Melting solid \ce{CO2} requires overcoming
  \choices
    \choice{covalent C=O bonds}
    \correct{London dispersion forces between molecules}
    \choice{ionic attractions}
    \choice{hydrogen bonds}
\end{mcq}
\why{Molecular solid: only intermolecular forces are disrupted.}
\distractor{A}{would decompose the compound, not melt it}

\begin{mcq}
  During the melting of a pure substance at constant pressure, the added
  heat causes
  \choices
    \choice{the temperature to rise steadily}
    \correct{intermolecular forces to be overcome at constant temperature}
    \choice{covalent bonds to break}
    \choice{the substance to become less dense in all cases}
\end{mcq}
\why{Phase-change plateau: energy goes into separating particles.}

\begin{mcq}
  A liquid boils when its vapor pressure
  \choices
    \choice{reaches zero}
    \choice{exceeds its critical pressure}
    \correct{equals the external pressure}
    \choice{equals its surface tension}
\end{mcq}
\why{Which is why water boils below \SI{100}{\celsius} at altitude.}

% ---- 3.4 gas laws (6) ------------------------------------------------------
\begin{mcq}
  A gas occupies \SI{4.0}{\liter} at \SI{300}{\kelvin}. At constant
  pressure, its volume at \SI{600}{\kelvin} is
  \choices
    \choice{\SI{2.0}{\liter}}
    \correct{\SI{8.0}{\liter}}
    \choice{\SI{4.0}{\liter}}
    \choice{\SI{16}{\liter}}
\end{mcq}
\why{Charles's law: $V \propto T$ in kelvin.}

\begin{mcq}
  Which quantity must be expressed in kelvin in $PV = nRT$?
  \choices
    \choice{pressure}
    \choice{volume}
    \correct{temperature}
    \choice{moles}
\end{mcq}
\why{Absolute temperature is required; Celsius has an arbitrary zero.}

\begin{mcq}
  A \SI{1.0}{\liter} flask contains \SI{0.20}{\mole} \ce{He} and
  \SI{0.30}{\mole} \ce{Ne}. The mole fraction of helium is
  \choices
    \choice{0.20}
    \correct{0.40}
    \choice{0.50}
    \choice{0.67}
\end{mcq}
\why{$0.20/0.50 = 0.40$.}

\begin{mcq}
  In a mixture of gases, the partial pressure of one component depends on
  \choices
    \choice{the identity of the other gases}
    \choice{the molar mass of that gas}
    \correct{its mole fraction and the total pressure}
    \choice{the sizes of the gas molecules}
\end{mcq}
\why{$P_A = \chi_A P_{\text{total}}$ --- identity is irrelevant for an ideal
gas.}

\begin{mcq}
  Two identical flasks at the same temperature and pressure contain \ce{H2}
  and \ce{CO2}. Which statement is correct?
  \choices
    \correct{They contain the same number of molecules.}
    \choice{They have the same mass.}
    \choice{The \ce{H2} flask has more molecules.}
    \choice{The \ce{CO2} flask has higher pressure.}
\end{mcq}
\why{Equal $P$, $V$, $T$ means equal $n$ (Avogadro's principle).}

% ---- 3.5 KMT (5) -----------------------------------------------------------
\begin{mcq}
  According to kinetic molecular theory, the average kinetic energy of gas
  particles is directly proportional to
  \choices
    \choice{pressure}
    \choice{molar mass}
    \correct{absolute temperature}
    \choice{container volume}
\end{mcq}
\why{$KE_{\text{avg}} = \frac32 k_BT$.}

\begin{mcq}
  At the same temperature, which gas has the highest average molecular
  speed?
  \choices
    \correct{\ce{H2}}
    \choice{\ce{O2}}
    \choice{\ce{CO2}}
    \choice{\ce{Kr}}
\end{mcq}
\why{$u_{\text{rms}} = \sqrt{3RT/M}$ --- lightest moves fastest.}

\begin{mcq}
  As the temperature of a gas increases, its Maxwell--Boltzmann distribution
  \choices
    \choice{narrows and shifts left}
    \correct{broadens, flattens, and shifts right}
    \choice{keeps the same shape}
    \choice{increases in total area}
\end{mcq}
\why{Total area is fixed by the number of molecules; the spread grows.}
\distractor{D}{area equals total molecules, which does not change}

\begin{mcq}
  Which KMT assumption explains why gas pressure increases when the
  container volume decreases at constant temperature?
  \choices
    \correct{particles are in constant motion and collide with the walls}
    \choice{collisions are perfectly elastic}
    \choice{particles have negligible volume}
    \choice{particles exert no attractive forces}
\end{mcq}
\why{Smaller volume means more wall collisions per unit area per second.}

\begin{mcq}
  Helium and argon are held at the same temperature. Compared with argon,
  helium has
  \choices
    \choice{greater average kinetic energy and greater speed}
    \correct{equal average kinetic energy and greater speed}
    \choice{equal average kinetic energy and equal speed}
    \choice{lower average kinetic energy and greater speed}
\end{mcq}
\why{KE depends only on $T$; speed depends on molar mass.}

% ---- 3.6 real gases (3) ----------------------------------------------------
\begin{mcq}
  Real gases behave most ideally at
  \choices
    \choice{low temperature and high pressure}
    \correct{high temperature and low pressure}
    \choice{low temperature and low pressure}
    \choice{high temperature and high pressure}
\end{mcq}
\why{Far apart (volume negligible) and fast-moving (attractions negligible).}

\begin{mcq}
  At moderate pressure and low temperature, $PV/nRT$ for a real gas dips
  below 1 primarily because
  \choices
    \correct{intermolecular attractions reduce the force of wall collisions}
    \choice{the particles occupy significant volume}
    \choice{the gas begins to decompose}
    \choice{the temperature is not absolute}
\end{mcq}
\why{Attractions lower measured $P$ relative to ideal.}
\distractor{B}{that effect pushes the ratio ABOVE 1 at high pressure}

\begin{mcq}
  Which gas is expected to deviate most from ideal behavior?
  \choices
    \choice{\ce{He}}
    \choice{\ce{N2}}
    \correct{\ce{H2O}}
    \choice{\ce{Ne}}
\end{mcq}
\why{Strong hydrogen bonding gives large attractive corrections.}

% ---- 3.7-3.8 solutions (3) -------------------------------------------------
\begin{mcq}
  What is the concentration of chloride ion in \SI{0.20}{\Molar} \ce{AlCl3}?
  \choices
    \choice{\SI{0.20}{\Molar}}
    \choice{\SI{0.40}{\Molar}}
    \correct{\SI{0.60}{\Molar}}
    \choice{\SI{0.067}{\Molar}}
\end{mcq}
\why{Three chlorides per formula unit.}

\begin{mcq}
  Diluting \SI{50.0}{\milli\liter} of \SI{4.00}{\Molar} solution to
  \SI{200.}{\milli\liter} gives a concentration of
  \choices
    \choice{\SI{0.500}{\Molar}}
    \correct{\SI{1.00}{\Molar}}
    \choice{\SI{2.00}{\Molar}}
    \choice{\SI{16.0}{\Molar}}
\end{mcq}
\why{$M_1V_1 = M_2V_2$: $(4.00)(50.0)/200. = 1.00$.}

\begin{mcq}
  Which particulate drawing correctly represents aqueous \ce{K2SO4}?
  \choices
    \choice{intact \ce{K2SO4} units dispersed in water}
    \correct{twice as many \ce{K+} particles as \ce{SO4^2-} particles, all
      separated}
    \choice{equal numbers of separated \ce{K+} and \ce{SO4^2-}}
    \choice{\ce{K+} and \ce{S^2-} and \ce{O^2-} all separated}
\end{mcq}
\why{Complete dissociation in the 2:1 ratio; the sulfate ion stays intact.}
\distractor{D}{polyatomic ions do not break apart on dissolving}

% ---- 3.9-3.10 solubility/separation (3) ------------------------------------
\begin{mcq}
  Which is most soluble in water?
  \choices
    \choice{\ce{CCl4}}
    \choice{\ce{C6H14}}
    \correct{\ce{CH3OH}}
    \choice{\ce{I2}}
\end{mcq}
\why{Methanol hydrogen bonds with water; the others are nonpolar.}

\begin{mcq}
  Distillation separates a mixture based on differences in
  \choices
    \choice{particle size}
    \correct{boiling point}
    \choice{color}
    \choice{molar mass}
\end{mcq}
\why{The more volatile component vaporizes first.}

\begin{mcq}
  In paper chromatography with a polar stationary phase, a component that
  travels only a short distance is best described as
  \choices
    \correct{strongly attracted to the stationary phase}
    \choice{very soluble in the mobile phase}
    \choice{the lightest component}
    \choice{nonpolar}
\end{mcq}
\why{Strong stationary-phase attraction retards movement.}

% ---- 3.11-3.13 spectroscopy (3) --------------------------------------------
\begin{mcq}
  Which type of radiation has the highest energy per photon?
  \choices
    \choice{microwave}
    \choice{infrared}
    \choice{visible}
    \correct{ultraviolet}
\end{mcq}
\why{$E = hc/\lambda$ --- shortest wavelength has highest energy.}

\begin{mcq}
  Infrared radiation is absorbed by molecules primarily to cause
  \choices
    \choice{electronic transitions}
    \correct{bond vibrations}
    \choice{molecular rotation}
    \choice{nuclear transitions}
\end{mcq}
\why{IR energies match vibrational spacings.}

\begin{mcq}
  In a Beer--Lambert calibration plot of absorbance versus concentration,
  the slope of the line equals
  \choices
    \choice{the concentration of the unknown}
    \correct{$\epsilon b$}
    \choice{the absorbance of the unknown}
    \choice{$1/\epsilon b$}
\end{mcq}
\why{$A = \epsilon bc$ is linear in $c$ with slope $\epsilon b$.}
\distractor{A}{a very common slip --- you must DIVIDE $A$ by the slope}

\examsection{II}{Free Response}{2 questions \textbullet\ 31 minutes}{\calculatorok}

% ---- Long FRQ --------------------------------------------------------------
\frq{long}{10}{CED 3.1--3.6 synthesis \textbullet\ SP 1, 4, 6}

Answer all parts. The table gives data for three substances.

\renewcommand{\arraystretch}{1.3}
\begin{center}
\begin{tabular}{@{}lccc@{}}
\toprule
Substance & Molar mass (\si{\gram\per\mole}) & Boiling point (\si{\celsius}) & Polarity\\
\midrule
\ce{CH4}   & 16.0 & $-164$ & nonpolar\\
\ce{CH3OH} & 32.0 & $65$   & polar\\
\ce{CH3SH} & 48.1 & $6$    & polar\\
\bottomrule
\end{tabular}
\end{center}

\begin{parts}
  \item Identify all intermolecular forces present in each of the three
        substances. \answerlines[3]{\ce{CH4}: LDF only. \ce{CH3OH}: LDF,
        dipole--dipole, and hydrogen bonding (H on O). \ce{CH3SH}: LDF and
        dipole--dipole (no H-bonding --- S is not N, O, or F).}
  \item \ce{CH3SH} has the largest molar mass yet boils \SI{59}{\celsius}
        below \ce{CH3OH}. Explain. \answerlines[3]{\ce{CH3OH} hydrogen
        bonds, which is far stronger than \ce{CH3SH}'s dipole--dipole
        attraction. That difference outweighs \ce{CH3SH}'s advantage in
        dispersion forces from its larger electron cloud.}
  \item Predict which of the three has the highest vapor pressure at
        \SI{25}{\celsius} and justify. \answerlines[2]{\ce{CH4} --- weakest
        IMFs (dispersion only) and by far the lowest boiling point, so the
        greatest fraction of molecules escapes the liquid.}
  \item A \SI{2.00}{\liter} vessel holds \SI{0.150}{\mole} of \ce{CH4} gas
        at \SI{27}{\celsius}. Calculate the pressure in atm.
        \workspace[2.4cm]
        \keyonly{\begin{apcanswer}$T = \SI{300.}{\kelvin}$;
        $P = nRT/V = (0.150)(0.08206)(300.)/2.00 =
        \SI{1.85}{\atm}$\end{apcanswer}}
  \item The vessel is cooled to \SI{-100}{\celsius} and the measured
        pressure is found to be \emph{lower} than $PV = nRT$ predicts.
        Explain in terms of the failed assumption.
        \answerlines[3]{At low temperature the molecules move slowly enough
        that intermolecular attractions become significant. Attractions pull
        molecules away from the walls, softening collisions, so the measured
        pressure falls below the ideal prediction ($PV/nRT < 1$).}
\end{parts}

\begin{rubric}
  \rpoint{2}{(a) 1 pt for \ce{CH4} and \ce{CH3OH} correct; 1 pt for
    \ce{CH3SH} WITHOUT hydrogen bonding --- listing H-bonding for S earns 0
    for this point}
  \rpoint{2}{(b) 1 pt names hydrogen bonding in \ce{CH3OH}; 1 pt states it
    outweighs the dispersion advantage of \ce{CH3SH}}
  \rpoint{1}{(c) \ce{CH4} with a weak-IMF justification}
  \rpoint{2}{(d) 1 pt converts to \SI{300.}{\kelvin}; 1 pt
    \SI{1.85}{\atm} with units}
  \rpoint{3}{(e) 1 pt identifies attractions as the failed assumption; 1 pt
    links attractions to softened wall collisions; 1 pt states measured $P$
    is below ideal. ``Because it is a real gas'' earns 0}
\end{rubric}

% ---- Short FRQ -------------------------------------------------------------
\frq{short}{4}{CED 3.13 Beer--Lambert \textbullet\ SP 4, 5}

A student measures the absorbance of four standard solutions of a blue dye
at \SI{630}{\nano\meter} in a \SI{1.00}{\centi\meter} cell and obtains a
straight line through the origin with slope
\SI{45.0}{\per\Molar}.

\begin{parts}
  \item Determine the molar absorptivity $\epsilon$, with units.
        \answerlines[2]{$\epsilon = \text{slope}/b = 45.0/1.00 =
        \SI{45.0}{\per\Molar\per\centi\meter}$.}
  \item An unknown solution gives $A = 0.630$. Calculate its
        concentration. \workspace[2cm]
        \keyonly{\begin{apcanswer}$c = A/(\epsilon b) = 0.630/45.0 =
        \SI{0.0140}{\Molar}$\end{apcanswer}}
  \item The student repeats the measurement at \SI{500}{\nano\meter} and
        obtains a much smaller absorbance. Explain why the original
        calibration line cannot be used at this new wavelength.
        \answerlines[2]{Molar absorptivity is wavelength-dependent, so
        $\epsilon$ --- and therefore the slope --- differs at
        \SI{500}{\nano\meter}; a new calibration is required.}
\end{parts}

\begin{rubric}
  \rpoint{1}{(a) \SI{45.0}{\per\Molar\per\centi\meter} with $b$ divided out}
  \rpoint{2}{(b) 1 pt divides $A$ by the slope (multiplying earns 0);
    1 pt \SI{0.0140}{\Molar} with units}
  \rpoint{1}{(c) states that $\epsilon$ depends on wavelength}
\end{rubric}

\examtotal

\end{document}
